\documentclass{article}

% Main-track anonymous submission.
\usepackage{neurips_2026}

\usepackage[utf8]{inputenc}
\usepackage[T1]{fontenc}
\usepackage{hyperref}
\usepackage{url}
\usepackage{booktabs}
\usepackage{amsfonts}
\usepackage{amsmath}
\usepackage{amssymb}
\usepackage{nicefrac}
\usepackage{microtype}
\usepackage{xcolor}

% Keep this file as a TEMPLATE, not a full draft.
% For project-specific wording and optional drafted prose, see:
%   NEURIPS_PAPER_START.md

\title{[Working Title]}

\author{%
 Anonymous Author(s) \\
 Anonymous Institution \\
 Anonymous Address \\
 \texttt{anonymous@email.com}
}

\begin{document}

\maketitle

\begin{abstract}
% 1 sentence: what current evaluations miss
% 1 sentence: what benchmark / framing is proposed
% 1 sentence: what environments / evidence support it
% 1 sentence: why it matters
[Abstract goes here.]
\end{abstract}

\section{Introduction}
% Fill only the core logic:
% - why outcome-only evaluation is insufficient
% - what hidden-mechanic discovery means in this project
% - the benchmark framing (World / Goal / Mechanics / Feedback)
% - the paper's contribution and roadmap
[Introduction goes here.]

\paragraph{Contributions.}
\begin{itemize}
  \item [Contribution 1: benchmark framing]
  \item [Contribution 2: environment family]
  \item [Contribution 3: metrics / epistemic traces]
  \item [Contribution 4: empirical finding]
\end{itemize}

\section{Related Work}
% Organize only around the buckets that matter:
% - POMDPs / partial observability
% - interactive agent benchmarks
% - latent-dynamics / rule-discovery / exploration
% - behavior-trace / process-oriented evaluation
[Related work goes here.]

\section{Method}

\subsection{Problem Setup}
% Define:
% - environment family
% - partial specification
% - World / Goal / Mechanics / Feedback axes
% - outcome metrics vs discovery metrics vs epistemic traces
[Problem setup goes here.]

\subsection{Benchmark Family}
% Explain how environments fit the family.
% Keep this high-level in the template.
% Candidate environments currently in scope:
% - Zendo (historical / earlier benchmark work)
% - KA59 (clean mechanics-discovery slice)
% - BP35 (interaction-heavy planning case)
% - LS20 (cross-environment support)
[Benchmark family description goes here.]

\subsection{Metrics}
% Suggested metric groups:
% - Outcome
% - Discovery
% - Epistemic signal
[Metric definitions go here.]

\section{Experiments}

\subsection{Experimental Setup}
% Fill in:
% - environments actually used in the current paper version
% - agents compared
% - conditions / axis manipulations
% - evaluation protocol
% - number of runs / seeds
[Experimental setup goes here.]

\subsection{Main Results}
% Recommended table logic:
% Environment | Condition | Agent | Outcome | Discovery | Epistemic signal

\begin{table}[t]
 \caption{Primary results table. Replace with real conditions, agents, and metrics.}
 \label{tab:results}
 \centering
 \begin{tabular}{llllll}
 \toprule
 Environment & Condition & Agent & Outcome & Discovery & Epistemic \\
 \midrule
 [env] & [cond] & [agent] & [x] & [x] & [x] \\
 \bottomrule
 \end{tabular}
\end{table}

[Main results text goes here.]

\subsection{Ablations and Analysis}
% Candidate analyses:
% - observability ablations
% - mechanics contrastive pairs
% - outcome/discovery divergence
% - epistemic-trace case studies
[Ablations and analysis go here.]

\section{Limitations}
% Be explicit about:
% - current environment count / maturity
% - agent strength
% - partial axis coverage
% - hand-designed epistemic metrics
% - benchmark-family incompleteness
[Limitations go here.]

\section{Conclusion}
% Recap the benchmark claim, main evidence, and next step.
[Conclusion goes here.]

\section*{References}
% Add bibliography or switch to BibTeX.
[References go here.]

\appendix

\section{Appendix}
% Environment specs, extra tables, implementation details, case studies.
[Appendix goes here.]

\newpage
\input{checklist.tex}

\end{document}
